\section{Consolidação da Fase 2 (M9)}
\label{sec:m9}

A milestone~M9 fecha os \emph{follow-ups} da Fase 2: leva a agregação \emph{max} dos testes
sintéticos para os dados reais e decide o \emph{default}; decide a Etapa 2 do multivariado
(OHLCV completo \emph{vs.} Close+Volume); e valida o gargalo do modelo multivariado.

\subsection{Adoção de \emph{max} como \emph{default} (ADR-0009)}
\label{sec:m9-max}
O notebook \texttt{10\_max\_default\_decision} reexecutou a detecção no \textbf{teste real
(2020--2024)} sob cada agregação. A fração média de janelas marcadas (quatro ativos) foi
$0{,}108$ (\texttt{mean}), $\mathbf{0{,}103}$ (\texttt{max}) e $0{,}112$ (\texttt{percentile})
com limiar estático, e $0{,}064$/$0{,}080$/$0{,}048$ com o dinâmico. Ou seja, \emph{max}
\textbf{não infla} a taxa de marcação no teste real --- fica praticamente igual a \emph{mean},
próxima da taxa-base do p95. Somado ao ganho de \emph{recall}/\emph{precision} na injeção
sintética (Seção~\ref{sec:m8-pooling-result}), a troca é segura: \texttt{config.yaml} passa a
\texttt{aggregation: max} por \emph{default}. Os notebooks de M4--M6 preservam suas saídas
originais (sob \texttt{mean}) como registro histórico.

\subsection{OHLCV completo \emph{vs.} Close+Volume (ADR-0011, Etapa 2)}
\label{sec:m9-ohlcv}
O notebook \texttt{11\_ohlcv\_full} treinou $(30,5)$ e $(30,2)$ nos quatro ativos
(Tabela~\ref{tab:m9-ohlcv}). O OHLCV completo \textbf{piora} o \texttt{val\_loss} em dois dos
quatro ativos e a atribuição do pico de volume é idêntica nas duas configurações. \textbf{Decisão:
permanecer em Close+Volume $(30,2)$}; a Etapa 2 não compensa o custo de reconstruir três canais
adicionais. Ressalva honesta: os \emph{log-retornos} de O/H/L/C \emph{não} são $\sim$99\%
colineares (correlação observada $0{,}39$--$0{,}75$; o $\sim$99\% referia-se aos níveis de
preço) --- mas, mesmo não-redundantes, não melhoraram a detecção.

\begin{table}[h]
\centering
\caption{\texttt{val\_loss} de treino: Close+Volume $(30,2)$ vs.\ OHLCV $(30,5)$ (ADR-0011, M9).}
\label{tab:m9-ohlcv}
\begin{tabular}{lcc}
\toprule
\textbf{Ativo} & \textbf{Close+Volume $(30,2)$} & \textbf{OHLCV $(30,5)$} \\
\midrule
PETR4.SA & 0,00214 & 0,00197 \\
VALE3.SA & \textbf{0,00232} & 0,00273 \\
AMER3.SA & \textbf{0,00254} & 0,00355 \\
ITUB4.SA & 0,00237 & 0,00240 \\
\bottomrule
\end{tabular}
\end{table}

\subsection{\texttt{latent\_dim} do modelo multivariado (ADR-0010/0011)}
\label{sec:m9-latent}
A malha walk-forward, generalizada para entrada por coluna
(\texttt{walk\_forward\_splits\_multivariate}), avaliou o gargalo do modelo Close+Volume:
\texttt{val\_loss} médio $0{,}01415$ ($8$), $0{,}01427$ ($16$), $0{,}01433$ ($32$) --- de novo
\textbf{insensível}, com leve vantagem para $8$. Mantém-se $16$ por consistência com o
univariado; a diferença é desprezível. A insensibilidade do gargalo, agora observada em ambas as
representações, é evidência consistente de robustez da reconstrução.
